% ======================================================================
% capitulos/capitulo2.tex
% ----------------------------------------------------------------------
% ENTREGABLE 2 — BUSINESS CASE (Caso de Negocio)
% Empresa caso de estudio: WISP Ingeniería en Redes y Telecomunicaciones
% Contenido trasladado desde Entregable2_BusinessCase_GTI_WISP.docx
% ======================================================================
\chapter{Business Case — Caso de Negocio del Proyecto}
\label{chap:businesscase}

% ----------------------------------------------------------------------
% Descripción del producto (introducción del capítulo)
% ----------------------------------------------------------------------
El presente documento constituye el Business Case (Caso de Negocio) del Proyecto
de Evaluación de Fin de Carrera para la empresa WISP Ingeniería en Redes y
Telecomunicaciones. Es un documento formal que justifica la viabilidad y
conveniencia del proyecto de tecnologías de información, presentando el problema
identificado, los objetivos del proyecto, las alternativas evaluadas y la solución
seleccionada, sustentando su aporte de valor.

El análisis incluye la evaluación de beneficios tangibles e intangibles, la
estimación de la inversión inicial, los costos operativos y de mantenimiento, el
retorno esperado mediante indicadores financieros (ROI, Payback, VAN y TIR) y una
evaluación preliminar de riesgos, demostrando que la iniciativa contribuye
estratégicamente al desempeño organizacional de WISP.

\begin{upeubox}{Continuidad con el Entregable~1 (Diagnóstico)}
Este Business Case toma como insumo el diagnóstico organizacional previamente
desarrollado (contexto, cadena de valor, FODA, cinco fuerzas de Porter y problema
identificado). Mientras el diagnóstico respondió \emph{«¿cuál es el problema y por
qué importa?»}, este documento responde \emph{«¿vale la pena resolverlo, cómo y con
qué retorno?»}.
\end{upeubox}

% ======================================================================
\section{Justificación del Proyecto}
\label{sec:justificacion}

\subsection{Problema que se resuelve}

Conforme al diagnóstico realizado, WISP gestiona la operación de aproximadamente
455 clientes de internet por fibra óptica distribuidos en 13 zonas geográficamente
dispersas, mediante un sistema heredado deficiente (PHP) y procesos manuales
apoyados en WhatsApp y registros en papel. Esta forma de operar genera pérdida y
dispersión de información, ausencia de trazabilidad en la captación e instalación,
gestión manual y repetitiva de la red en cada MikroTik, y una excesiva dependencia
de una sola persona.

El proyecto resuelve este problema mediante el desarrollo e implementación de un
\textbf{Sistema de Gestión Integral para ISP de Fibra Óptica Multi-Zona}, compuesto
por un sistema web de administración, una aplicación móvil con roles de técnico y
vendedor, y la integración directa con los 13 MikroTiks vía VPN (WireGuard) y la API
de RouterOS, junto con la automatización de notificaciones por WhatsApp.

\subsection{Objetivos del proyecto (SMART)}

\begin{upeubox}{Fundamento metodológico: objetivos SMART}
Los objetivos se formulan siguiendo el criterio SMART propuesto por George~T. Doran
(1981), según el cual todo objetivo debe ser Específico (\textit{Specific}), Medible
(\textit{Measurable}), Alcanzable (\textit{Achievable}), Relevante (\textit{Relevant})
y Temporal (\textit{Time-bound}).
\end{upeubox}

\begin{upeubox}{Objetivo general}
Desarrollar e implementar, en un plazo de 5 meses, un sistema de gestión integral
(web $+$ móvil) que centralice y automatice los procesos de captación, instalación,
activación y control de clientes de WISP en sus 13 zonas, reemplazando el uso de
WhatsApp y el sistema heredado, y habilitando la escalabilidad de la operación.
\end{upeubox}

Los objetivos específicos, formulados bajo el criterio SMART, se presentan a
continuación junto con su métrica de verificación:

\begin{center}
\captionof{table}{Objetivos específicos del proyecto (SMART).}
\label{tab:objetivos-smart}
\begin{tabularx}{\textwidth}{C{1.1cm} X X}
\toprule
\textbf{ID} & \textbf{Objetivo específico} & \textbf{Métrica de verificación} \\
\midrule
OP-1 &
Digitalizar el 100\% del proceso de captación mediante una app móvil para
vendedores, eliminando el registro en papel/WhatsApp, al finalizar el proyecto. &
\% de prospectos registrados por app vs.\ total \\
OP-2 &
Implementar el registro estructurado del 100\% de las instalaciones técnicas
(fotos de caja NAP, borne, equipos, fachada) vinculado a cada cliente. &
\% de instalaciones con evidencia completa en el sistema \\
OP-3 &
Automatizar la gestión de usuarios PPPoE (alta, corte, reactivación) sobre los 13
MikroTiks desde el sistema central vía API de RouterOS. &
N° de MikroTiks integrados / 13 \\
OP-4 &
Centralizar en una única base de datos la información de los $\approx$455 clientes,
migrando los datos dispersos del sistema heredado y los MikroTiks. &
\% de clientes migrados y consolidados \\
OP-5 &
Reducir el tiempo administrativo de gestión diaria en al menos un 50\% respecto a la
operación manual actual, a los 3 meses de la puesta en marcha. &
Horas/día de gestión antes vs.\ después \\
\bottomrule
\end{tabularx}
\end{center}

\subsection{Beneficios esperados}

La ejecución del proyecto generará beneficios en tres niveles: operativo,
estratégico y para el cliente final.

\begin{itemize}[leftmargin=1.4em]
  \item \textbf{Operativos:} eliminación del retrabajo manual, reducción de errores,
  trazabilidad completa del ciclo de vida del cliente y automatización de la gestión
  de red.
  \item \textbf{Estratégicos:} habilitación de la escalabilidad para expandirse a
  nuevos distritos, reducción de la dependencia de una sola persona y disponibilidad
  de información para la toma de decisiones.
  \item \textbf{Para el cliente:} activación más rápida del servicio, atención más
  ágil de averías y comunicación más ordenada.
\end{itemize}

% ======================================================================
\section{Análisis de Alternativas}
\label{sec:alternativas}

\subsection{Alternativas tecnológicas consideradas}

Se evaluaron cuatro alternativas para resolver el problema identificado, desde no
intervenir hasta el desarrollo de una solución propia a medida:

\begin{center}
\captionof{table}{Alternativas tecnológicas evaluadas.}
\label{tab:alternativas}
\begin{tabularx}{\textwidth}{L{2.2cm} X X X}
\toprule
\textbf{Alternativa} & \textbf{Descripción} & \textbf{Ventaja} & \textbf{Desventaja} \\
\midrule
A1. Statu quo &
Mantener el sistema heredado (PHP) y la operación por WhatsApp. &
Costo cero inmediato. &
No resuelve nada; el problema se agrava con el crecimiento. \\
A2. Software ISP comercial &
Adquirir una plataforma existente (ej.\ WISPHub, Mikrowisp, SmartOLT). &
Rápido de adoptar; ya probado. &
Costo mensual por suscripción; poca personalización; dependencia del proveedor;
puede no cubrir apps de campo a medida. \\
A3. Desarrollo externo &
Contratar una empresa de software para desarrollo a medida. &
Solución personalizada. &
Costo elevado (S/~25\,000--40\,000); plazos largos; dependencia del proveedor para
mantenimiento. \\
A4. Desarrollo propio &
Desarrollar la solución a medida con el equipo del proyecto (4 integrantes). &
Personalización total; sin costo de desarrollo para WISP; conocimiento interno. &
Requiere tiempo del equipo; curva de aprendizaje. \\
\bottomrule
\end{tabularx}
\end{center}

\subsection{Criterios de evaluación}

\begin{upeubox}{Fundamento metodológico: modelo de puntuación ponderada}
La comparación se realiza mediante un modelo de puntuación ponderada
(\textit{weighted scoring model}), técnica de análisis multicriterio de decisiones.
A cada criterio se le asigna un peso relativo según su importancia; cada alternativa
se puntúa de 1 a 5 en cada criterio; el puntaje ponderado total determina la mejor
opción.
\end{upeubox}

Se definieron seis criterios de evaluación con sus respectivos pesos:

\begin{center}
\captionof{table}{Criterios de evaluación y pesos asignados.}
\label{tab:criterios}
\begin{tabularx}{\textwidth}{L{3.6cm} C{1.2cm} X}
\toprule
\textbf{Criterio} & \textbf{Peso} & \textbf{Justificación del peso} \\
\midrule
Ajuste a la necesidad & 25\% &
Qué tanto la solución cubre los 3 flujos (captación, instalación, gestión) y la
integración con MikroTik. \\
Costo total (inversión $+$ operación) & 20\% &
Impacto económico sobre una MYPE con recursos limitados. \\
Personalización & 20\% &
Capacidad de adaptarse a la operación específica de WISP (apps de campo, cobertura). \\
Tiempo de implementación & 15\% &
Rapidez en poner la solución en producción. \\
Escalabilidad & 12\% &
Capacidad de crecer con nuevas zonas y clientes. \\
Independencia del proveedor & 8\% &
Menor dependencia externa para operar y mantener el sistema. \\
\bottomrule
\end{tabularx}
\end{center}

\subsection{Matriz comparativa (scoring ponderado)}

Cada alternativa se puntuó de 1 (muy deficiente) a 5 (excelente) en cada criterio.
La columna final muestra el puntaje ponderado total sobre un máximo de 5.00:

\begin{center}
\captionof{table}{Matriz comparativa de alternativas (scoring ponderado).}
\label{tab:scoring}
{\small
\begin{tabularx}{\textwidth}{X C{1.05cm} C{1.05cm} C{1.05cm} C{1.05cm} C{1.05cm} C{1.05cm} C{1.1cm}}
\toprule
\textbf{Alternativa} & \textbf{Ajuste (25\%)} & \textbf{Costo (20\%)} &
\textbf{Person. (20\%)} & \textbf{Tiempo (15\%)} & \textbf{Escala (12\%)} &
\textbf{Indep. (8\%)} & \textbf{TOTAL} \\
\midrule
A1. Statu quo            & 1 & 5 & 1 & 5 & 1 & 3 & 2.56 \\
A2. Software comercial   & 3 & 3 & 2 & 5 & 3 & 2 & 3.02 \\
A3. Desarrollo externo   & 5 & 2 & 5 & 2 & 4 & 2 & 3.59 \\
\rowcolor{upeublue!10}
\textbf{A4. Desarrollo propio} $\star$ & \textbf{5} & \textbf{5} & \textbf{5} &
\textbf{3} & \textbf{5} & \textbf{5} & \textbf{4.70} \\
\bottomrule
\end{tabularx}
}
\end{center}

\begin{upeubox}{Alternativa seleccionada: A4 — Desarrollo propio}
Con un puntaje ponderado de \textbf{4.70/5.00}, la alternativa de desarrollo propio
es la de mayor valor. Supera a las demás por su ajuste total a la necesidad, su
personalización completa (apps de campo e integración con MikroTik a medida), su
nulo costo de desarrollo para WISP y su alta escalabilidad e independencia. Su única
desventaja ---el tiempo del equipo--- es asumible en el marco del proyecto de fin de
carrera.
\end{upeubox}

% ======================================================================
\section{Evaluación de Beneficios}
\label{sec:beneficios}

Los beneficios del proyecto se clasifican en cuantificables (traducibles a un valor
monetario) y cualitativos (de valor estratégico no monetario directo). Las
estimaciones son conservadoras y se basan en los datos operativos de WISP.

\subsection{Beneficios cuantificables (anuales)}

Sobre una base estimada de $\approx$455 clientes activos y un ticket promedio de
S/~55 mensuales (ingreso anual aproximado de S/~300\,300), se estiman los siguientes
beneficios económicos anuales:

\begin{center}
\captionof{table}{Beneficios cuantificables anuales estimados.}
\label{tab:beneficios-cuant}
\begin{tabularx}{\textwidth}{L{4cm} X R{2cm}}
\toprule
\textbf{Beneficio cuantificable} & \textbf{Cálculo / supuesto} & \textbf{Valor anual (S/)} \\
\midrule
Ahorro de tiempo administrativo &
2.5 h/día liberadas $\times$ 26 días $\times$ S/~10/h $\times$ 12 meses. & 7\,800 \\
Recuperación por menor morosidad &
2\% del ingreso anual recuperado por cortes/reactivaciones oportunos y automatizados. & 6\,006 \\
Ingreso por clientes extra captados &
3 clientes nuevos extra/mes por mejor captación $\times$ S/~55 $\times$ $\sim$6 meses
activos promedio. & 11\,880 \\
Reducción de errores y pérdida de prospectos &
Menor pérdida de datos de campo, papel y retrabajo de coordinación. & 1\,800 \\
\midrule
\textbf{BENEFICIO BRUTO ANUAL} & \textbf{Suma de beneficios cuantificables} & \textbf{27\,486} \\
\bottomrule
\end{tabularx}
\end{center}

\subsection{Beneficios cualitativos}

\begin{itemize}[leftmargin=1.4em]
  \item \textbf{Trazabilidad total:} historial completo por cliente (instalación,
  caja NAP, borne, equipos), valioso para soporte y mantenimiento.
  \item \textbf{Reducción del riesgo operativo:} la operación deja de depender de una
  sola persona y de WhatsApp.
  \item \textbf{Escalabilidad:} capacidad de crecer a nuevos distritos sin aumentar
  proporcionalmente la carga administrativa.
  \item \textbf{Imagen y profesionalización:} una operación digitalizada mejora la
  percepción de la empresa frente a clientes y competidores.
  \item \textbf{Información para decisiones:} datos centralizados que permiten
  analizar zonas, morosidad y crecimiento.
\end{itemize}

\subsection{Indicadores de valor}

Los principales indicadores de valor (KPI) que evidenciarán el beneficio del proyecto
tras su implementación son:

\begin{center}
\captionof{table}{Indicadores de valor (KPI) del proyecto.}
\label{tab:kpi}
\begin{tabularx}{\textwidth}{L{4.2cm} X X}
\toprule
\textbf{Indicador de valor (KPI)} & \textbf{Situación actual} & \textbf{Meta con el sistema} \\
\midrule
Tiempo de gestión administrativa diaria & $\sim$4--5 h/día & $\le$ 2 h/día \\
Tiempo de activación de un cliente nuevo & Variable (días) & $<$ 24 h desde validación \\
Trazabilidad de instalaciones & 0\% (WhatsApp) & 100\% en el sistema \\
Prospectos registrados digitalmente & 0\% (papel/WhatsApp) & 100\% por app \\
Tasa de morosidad & Sin control preciso & Reducción $\ge$ 2\% \\
\bottomrule
\end{tabularx}
\end{center}

% ======================================================================
\section{Estimación de Costos}
\label{sec:costos}

La estimación de costos distingue entre la inversión inicial (única, al arranque),
los costos operativos y de mantenimiento (recurrentes) y el valor de mercado del
desarrollo (que WISP no paga, pues lo aporta el equipo del proyecto, pero que se
cuantifica para medir el ahorro y el retorno real).

\subsection{Inversión inicial (asumida por WISP)}

El modelo del proyecto establece que el equipo de 4 integrantes desarrolla e
implementa la solución, mientras que WISP asume únicamente la infraestructura de
producción y la puesta en marcha:

\begin{center}
\captionof{table}{Inversión inicial asumida por WISP.}
\label{tab:inversion}
\begin{tabularx}{\textwidth}{L{4.2cm} X R{1.8cm}}
\toprule
\textbf{Concepto} & \textbf{Detalle} & \textbf{Costo (S/)} \\
\midrule
VPS en la nube (setup $+$ 1.er mes) & Servidor cloud nuevo para producción. & 150 \\
Dominio web & Registro de dominio nuevo (.com/.pe), anual. & 60 \\
Certificado SSL / correo corporativo & Seguridad y correo institucional. & 90 \\
Migración de datos y puesta en marcha & Migración de clientes del sistema heredado
y MikroTiks. & 300 \\
Capacitación al personal & 2 sesiones para admin, técnicos y vendedores. & 200 \\
Contingencia / imprevistos & Reserva para imprevistos ($\sim$10\%). & 200 \\
\midrule
\textbf{TOTAL INVERSIÓN INICIAL} & \textbf{Desembolso real de WISP al arranque} & \textbf{1\,000} \\
\bottomrule
\end{tabularx}
\end{center}

\subsection{Costos operativos y de mantenimiento (anuales)}

\begin{center}
\captionof{table}{Costos operativos y de mantenimiento anuales.}
\label{tab:costos-op}
\begin{tabularx}{\textwidth}{L{4.2cm} X R{1.8cm}}
\toprule
\textbf{Concepto} & \textbf{Detalle} & \textbf{Costo anual (S/)} \\
\midrule
VPS mensual & Hosting del sistema en producción (S/~130 $\times$ 12). & 1\,560 \\
Renovación de dominio & Renovación anual del dominio. & 60 \\
Mantenimiento y soporte evolutivo & Ajustes, mejoras y soporte (S/~150 $\times$ 12). & 1\,800 \\
Respaldo / backup en nube & Copias de seguridad (S/~30 $\times$ 12). & 360 \\
\midrule
\textbf{TOTAL COSTO OPERATIVO ANUAL} & \textbf{$\approx$ S/~315 mensuales} & \textbf{3\,780} \\
\bottomrule
\end{tabularx}
\end{center}

\begin{upeubox}{Valor de mercado del desarrollo (ahorro para WISP)}
Si WISP hubiera contratado el desarrollo de este sistema (web $+$ app móvil $+$
integración con MikroTik) a una empresa de software, el costo de mercado se estima en
aproximadamente \textbf{S/~30\,000}. Al ser desarrollado por el equipo del proyecto de
fin de carrera, WISP no incurre en este costo, lo que representa un ahorro directo de
esa magnitud. Este valor se utiliza como referencia en el escenario riguroso de la
evaluación financiera.
\end{upeubox}

% ======================================================================
\section{Retorno Esperado — Evaluación Financiera}
\label{sec:retorno}

\begin{upeubox}{Fundamento metodológico: indicadores financieros}
La evaluación se realiza con los indicadores de finanzas corporativas: ROI (retorno
sobre la inversión), Período de Recuperación (Payback), VAN (Valor Actual Neto) y TIR
(Tasa Interna de Retorno). VAN y TIR incorporan el valor del dinero en el tiempo
mediante una tasa de descuento del 12\% anual, referencial para proyectos de TI en el
contexto peruano (ajustable por la empresa).
\end{upeubox}

El beneficio neto anual del proyecto se obtiene restando los costos operativos al
beneficio bruto anual: S/~27\,486 $-$ S/~3\,780 $=$ \textbf{S/~23\,706}. Se presentan
dos escenarios de evaluación:

\begin{itemize}[leftmargin=1.4em]
  \item \textbf{Escenario A (costo real para WISP):} considera solo el desembolso real
  de la empresa (S/~1\,000). Refleja el retorno efectivo que obtiene WISP.
  \item \textbf{Escenario B (riguroso, con valor de mercado):} considera la inversión
  como si se hubiera pagado el desarrollo a precio de mercado (S/~30\,000). Es el
  análisis financiero conservador y defendible.
\end{itemize}

\begin{center}
\captionof{table}{Indicadores financieros por escenario.}
\label{tab:financiero}
\begin{tabularx}{\textwidth}{L{2.8cm} C{2.3cm} C{2.5cm} X}
\toprule
\textbf{Indicador} & \textbf{Escenario A (costo real WISP)} &
\textbf{Escenario B (valor de mercado)} & \textbf{Interpretación} \\
\midrule
Inversión & S/~1\,000 & S/~30\,000 & Desembolso considerado. \\
Beneficio neto anual & S/~23\,706 & S/~23\,706 & Igual en ambos: el flujo es el mismo. \\
ROI (año 1) & 1\,876\% & $-$21\% & En B se recupera desde el año 2. \\
Payback & 0.6 meses & 15.2 meses & Recuperación rápida en ambos casos. \\
VAN (3 años, 12\%) & S/~55\,738 & S/~26\,938 & VAN $>0$ \textrightarrow{} proyecto viable. \\
TIR & $>$1\,000\% & 60\% & TIR $\gg$ 12\% \textrightarrow{} muy rentable. \\
\bottomrule
\end{tabularx}
\end{center}

\begin{upeubox}{Conclusión de la evaluación financiera}
En ambos escenarios el proyecto es financieramente viable: el VAN es positivo y la
TIR supera ampliamente la tasa de descuento del 12\%. Incluso bajo el escenario
riguroso B ---que asume el costo de mercado completo del desarrollo--- el proyecto se
recupera en poco más de 15 meses y genera un VAN de S/~26\,938 en tres años. Bajo el
escenario real A, el retorno para WISP es prácticamente inmediato. La inversión está
plenamente justificada.
\end{upeubox}

% ======================================================================
\section{Riesgos Iniciales}
\label{sec:riesgos}

\begin{upeubox}{Fundamento metodológico: enfoque PMBOK (PMI)}
Los riesgos se evalúan siguiendo el enfoque de la Guía PMBOK (PMI): cada riesgo se
valora por su Probabilidad y su Impacto, y el producto de ambos determina su nivel de
severidad (Bajo, Medio, Alto), priorizando así la atención y las respuestas.
\end{upeubox}

\subsection{Identificación de riesgos estratégicos}

Se identifican los principales riesgos que podrían afectar la ejecución y adopción
del proyecto:

\begin{center}
\captionof{table}{Identificación de riesgos estratégicos.}
\label{tab:riesgos-id}
\begin{tabularx}{\textwidth}{C{1cm} X L{2.8cm}}
\toprule
\textbf{ID} & \textbf{Riesgo} & \textbf{Categoría} \\
\midrule
R1 & Baja adopción del sistema por el personal (técnicos/vendedoras) acostumbrado a
WhatsApp. & Organizacional \\
R2 & Fallo o incompatibilidad en la integración con los MikroTik vía API en alguna
zona. & Técnico \\
R3 & Pérdida o inconsistencia de datos durante la migración desde el sistema
heredado. & Técnico / Datos \\
R4 & Conectividad inestable en zonas rurales que afecte el uso de las apps en campo. &
Operativo \\
R5 & Retrasos en el desarrollo por la carga académica del equipo. & Proyecto \\
R6 & Dependencia del VPS en la nube: caída del servicio afectaría la gestión central. &
Infraestructura \\
R7 & Brechas de seguridad / acceso no autorizado a los datos de clientes. & Seguridad \\
\bottomrule
\end{tabularx}
\end{center}

\subsection{Evaluación preliminar de impacto (matriz probabilidad $\times$ impacto)}

Cada riesgo se valora en una escala de 1 (muy bajo) a 5 (muy alto) para probabilidad
e impacto. La severidad es el producto de ambos: 1--6 Bajo, 8--12 Medio, 15--25 Alto.

\begin{center}
\captionof{table}{Matriz de evaluación de riesgos y estrategias de mitigación.}
\label{tab:riesgos-matriz}
{\small
\begin{tabularx}{\textwidth}{C{0.8cm} C{0.9cm} C{0.9cm} C{0.9cm} L{1.4cm} X}
\toprule
\textbf{ID} & \textbf{Prob. (1-5)} & \textbf{Imp. (1-5)} & \textbf{Sev.} &
\textbf{Nivel} & \textbf{Estrategia de respuesta (mitigación)} \\
\midrule
R1 & 4 & 4 & 16 & Alto &
Capacitación práctica, interfaz simple tipo app conocida, acompañamiento inicial y
apoyo del administrador. \\
R2 & 3 & 4 & 12 & Medio &
Pruebas de integración por zona antes del despliegue; los 13 MikroTik usan el mismo
modelo y RouterOS 7 (homogéneo). \\
R3 & 3 & 5 & 15 & Alto &
Respaldo previo, migración por etapas, validación de datos y conservación del sistema
heredado hasta confirmar consistencia. \\
R4 & 4 & 3 & 12 & Medio &
Apps con modo offline / sincronización diferida para registrar en campo aun sin señal. \\
R5 & 4 & 3 & 12 & Medio &
Metodología ágil con entregables incrementales; priorización de módulos críticos;
cronograma con holgura. \\
R6 & 2 & 4 & 8 & Medio &
Backups automáticos, proveedor cloud confiable y plan de recuperación; posibilidad de
respaldo local. \\
R7 & 3 & 5 & 15 & Alto &
Autenticación por roles, cifrado, VPN WireGuard para acceso a MikroTik y buenas
prácticas de seguridad. \\
\bottomrule
\end{tabularx}
}
\end{center}

\begin{upeubox}{Conclusión de la evaluación de riesgos}
Los riesgos de mayor severidad son la baja adopción del personal (R1), la pérdida de
datos en la migración (R3) y la seguridad de la información (R7). Todos cuentan con
estrategias de mitigación viables dentro del alcance del proyecto. Ninguno constituye
un impedimento crítico: el proyecto es ejecutable con un nivel de riesgo gestionable.
\end{upeubox}

% ======================================================================
\section{Conclusión del Business Case}
\label{sec:conclusion-bc}

El análisis desarrollado demuestra que el proyecto es conveniente, viable y
estratégicamente alineado con los objetivos de WISP:

\begin{itemize}[leftmargin=1.4em]
  \item \textbf{Justificado:} resuelve un problema real y crítico (operación manual,
  sin trazabilidad y dependiente de una persona), con objetivos SMART verificables.
  \item \textbf{Óptimo:} el análisis multicriterio selecciona el desarrollo propio
  (4.70/5.00) como la mejor alternativa frente a mantener lo actual, comprar software
  comercial o contratar desarrollo externo.
  \item \textbf{Rentable:} con VAN positivo y TIR muy superior al costo de capital en
  ambos escenarios, y un Payback de entre 0.6 y 15 meses.
  \item \textbf{Gestionable:} los riesgos identificados tienen respuestas de mitigación
  viables y ninguno es un impedimento crítico.
\end{itemize}

\begin{upeubox}{Recomendación final}
Se recomienda \textbf{aprobar y ejecutar el proyecto}. La inversión requerida por WISP
es mínima frente a los beneficios operativos, estratégicos y económicos esperados, y
la iniciativa habilita directamente la escalabilidad y competitividad de la empresa en
su mercado. Este Business Case constituye la base para el Plan de Gestión del Proyecto
(Entregable~3).
\end{upeubox}
